\documentclass[11pt,letterpaper]{article}
\usepackage[T1]{fontenc}
\usepackage{lmodern}
\usepackage[margin=0.9in,headheight=14pt]{geometry}
\usepackage{amsmath,amssymb,amsthm,mathtools}
\usepackage{booktabs,array}
\usepackage{xcolor,microtype,fancyhdr,hyperref}
\definecolor{ink}{HTML}{17354C}
\definecolor{pale}{HTML}{F2F5F7}
\hypersetup{colorlinks=true,linkcolor=ink,citecolor=ink,urlcolor=ink,
 pdftitle={IE-27: Continued analysis and verified extensions},pdfauthor={Sidney Holden}}
\pagestyle{fancy}\fancyhf{}
\lhead{\small\textcolor{ink}{IE-27 / Continued investigation}}
\rhead{\small 12 September 2026}
\cfoot{\small\thepage}
\setlength{\parindent}{0pt}\setlength{\parskip}{6pt}
\setlength{\emergencystretch}{2em}\clubpenalty=10000\widowpenalty=10000
\newtheorem{theorem}{Theorem}[section]
\newtheorem{proposition}[theorem]{Proposition}
\newtheorem{lemma}[theorem]{Lemma}
\newcommand{\R}{\mathbb R}\newcommand{\C}{\mathbb C}
\newcommand{\norm}[1]{\left\|#1\right\|_2}
\newcommand{\code}[1]{\texttt{\detokenize{#1}}}
\begin{document}
{\small\sffamily\bfseries\textcolor{ink}{MATHEMATICAL RESEARCH UPDATE}}

{\LARGE\bfseries\textcolor{ink}{IE-27: Continued analysis\\[3pt]and verified extensions}}

{\large The spectral disk for a Radau stage preconditioner}

\par\textbf{Sidney Holden}\par
{\small Center for Computational Biology, Flatiron Institute, Simons Foundation\par Prepared with substantial AI assistance}\par\vspace{5pt}
\fcolorbox{ink}{pale}{\parbox{0.948\linewidth}{
\textbf{Status: this archive is not a complete solution of IE-27.}
The unrestricted stage-number inequality is not proved here, and no
counterexample to the stated real-positive-shift conjecture is supplied.
The new results below do not remove that gap.}}

\section{Target and scope}
For the increasing Radau IIA nodes and their collocation matrix $A_q$, take
\[
 A_q^{-1}=L_qU_q,\qquad \operatorname{diag}(U_q)=\boldsymbol1,
 \qquad B_q=L_q^{-1},\quad N_q=U_q-I_q,\quad r_q=\norm{N_q}<1.
\]
The precise remaining target is
\begin{equation}\label{eq:target}
 \boxed{\rho\bigl((I_q+\mu B_q)^{-1}N_q\bigr)\le r_q}\qquad
 (q\ge2,\; r_q<1,\; \mu\in\mathbb R_{>0}).
\end{equation}
This is the small-matrix formulation of the requested disk enclosure for the
symmetric positive-definite spatial pencil. It is given on the problem page
and in the cited spectral reduction; the disk uses a \emph{radius}, not a
diameter \cite{problem,outrata}.

The prior nine-page report and its files are retained in
\code{prior_results/}. That report gives the spatial reduction, exact
Radau inverse formula, all-stage positive-pivot theorem, and exact two- and
three-stage arguments. Its 64 certificates cover $q=2,\ldots,64$ and $q=80$.

This continuation adds an all-stage small/large-shift theorem, a rank-two
weighted identity, a single-matrix sufficient condition for a stronger norm
bound, and rigorous certificates for $q=96$ and $q=128$. It also gives exact
counterexamples to two \emph{broader} assertions that must not be substituted
for the Radau conjecture.

\begin{center}
\begin{tabular}{@{}p{0.67\linewidth}p{0.25\linewidth}@{}}
\toprule
Claim & Status in this archive\\\midrule
Equation \eqref{eq:target} for every admissible stage number & Not established\\
All positive shifts for $q=2,\ldots,64,80,96,128$ & 66 exact certificates\\
Small and large positive shifts, every $q$ & Proved in Section 3\\
Complex right-half-plane shifts in place of real shifts & False; exact example\\
Generic matrices in place of Radau matrices & False; exact example\\\bottomrule
\end{tabular}
\end{center}

\newpage
\section{All-stage structure and a rank-two identity}
Let $0<c_1<\cdots<c_q=1$ be the specified nodes, and set
\[
 R(t)=\prod_{j=1}^q(t-c_j),\quad v_i=c_iR'(c_i),\quad
 V=\operatorname{diag}(v_i),\quad X=\operatorname{diag}(c_i).
\]
The inverse differentiation formula, proved in the retained report, is
\begin{equation}\label{eq:D}
 D_{ij}=\frac{v_i}{v_j(c_i-c_j)}\quad(i\ne j),\qquad
 D_{ii}=\frac1{2c_i}\ (i<q),\qquad D_{qq}=\frac{q^2+1}{2}.
\end{equation}
For completeness, this formula comes from the cardinal basis
$t\ell_j(t)/c_j$ of polynomials of degree at most $q$ vanishing at zero.
Its derivative evaluations give $A_q^{-1}$. The diagonal simplification
uses the equation
\[
 t(1-t)J''(t)+(1-3t)J'(t)+(q^2-1)J(t)=0,
 \qquad J(t)=P_{q-1}^{(1,0)}(2t-1),
\]
whose roots are the interior nodes. The retained report supplies the
coefficient-level derivation, so the certificate checker does not trust a
rounded Butcher matrix.

\begin{proposition}\label{prop:ranktwo}
For $C=V^{-1}DV$, $e=(1,\ldots,1)^T$, and the last coordinate vector $e_q$,
\begin{align}
 C+C^T&=\operatorname{diag}(1/c_1,\ldots,1/c_{q-1},q^2+1)\succ0,\label{eq:diag}\\
 XC+C^TX&=ee^T+q^2e_qe_q^T.\label{eq:ranktwo}
\end{align}
Every pivot of the exact no-pivot Crout factorization $D=LU$ is positive.
\end{proposition}
\begin{proof}
Off the diagonal, $C_{ij}=1/(c_i-c_j)$, so \eqref{eq:diag} follows by
cancellation. The off-diagonal entries of \eqref{eq:ranktwo} are
$(c_i-c_j)/(c_i-c_j)=1$. Its first $q-1$ diagonal entries are one, and
its last is $q^2+1$.

Every leading principal submatrix of $C$ has positive-definite symmetric
part. For each of its complex eigenpairs, the real part of the eigenvalue
is positive, as follows by taking the Hermitian inner product with its
eigenvector. Real eigenvalues are therefore positive, and conjugate pairs
have positive product. Every leading determinant is positive. The same
holds for $D$ by diagonal similarity. Successive ratios of these leading
determinants are exactly the pivots of the specified factorization.
\end{proof}

Consequently $B=L^{-1}$ exists, has positive diagonal entries, and
$I+\mu B$ is nonsingular for every real $\mu\ge0$. These facts hold
without first proving $r_q<1$.

Identity \eqref{eq:ranktwo} is useful additional structure, but it concerns
$D$, not the Euclidean symmetric part of $B$. Replacing that distinction
by an unproved accretivity assertion would invalidate a disk proof.

\newpage
\section{Two rigorously covered shift regimes at every stage}
This result uses only the triangular structure and the positive pivots;
it does not rely on a finite-stage computation.

\begin{theorem}\label{thm:shifts}
For every $q\ge2$, the requested enclosure holds whenever
\begin{equation}\label{eq:regimes}
 0<\mu\le\frac{1}{q\norm{B_q}}
 \qquad\text{or}\qquad
 \mu\ge2\norm{L_q}.
\end{equation}
The second regime in fact gives the stronger inequality
$\norm{(I+\mu B_q)^{-1}N_q}\le r_q$.
\end{theorem}
\begin{proof}
Suppress the subscript $q$ and write $r=\norm{N}$. Here $r>0$: for example,
$U_{12}=D_{12}/L_{11}\ne0$ by \eqref{eq:D}. Since $N$ is strictly upper
triangular, $N^q=0$.

Suppose $\lambda$ is an eigenvalue of $(I+\mu B)^{-1}N$ with
$|\lambda|>r$, and choose a nonzero eigenvector $x$. Then
\[
 (\lambda I-N)x=-\lambda\mu Bx,
 \qquad
 (\lambda I-N)^{-1}=\sum_{k=0}^{q-1}\frac{N^k}{\lambda^{k+1}}.
\]
Taking norms gives
\[
 1\le\mu\norm{B}\sum_{k=0}^{q-1}\left(\frac r{|\lambda|}\right)^k
 <q\mu\norm{B}.
\]
This is impossible in the first regime, including its endpoint.

For the second regime, use
\[
 (I+\mu B)^{-1}N=(L+\mu I)^{-1}LN.
\]
The Neumann-series estimate, valid for $\mu>\norm{L}$, yields
\[
 \norm{(I+\mu B)^{-1}N}
 \le \frac{\norm{L}r}{\mu-\norm{L}}\le r
 \quad\text{when }\mu\ge2\norm{L}.
\]
\end{proof}

A sharper small-shift version replaces $q$ by
\[
 \eta_q=\sum_{k=0}^{q-1}\frac{\norm{N_q^k}}{r_q^k}\le q.
\]
The same argument works because the corresponding resolvent sum is
strictly smaller than $\eta_q$ when $|\lambda|>r_q$.

\textbf{What this does not prove.} For an uncertified stage, a possible
violation is now confined to the finite interval
\[
 \frac{1}{q\norm{B_q}}<\mu<2\norm{L_q}.
\]
The theorem does not cover that interval. Both its endpoints still depend
on the stage, and their existence alone is not a uniform-stage solution.

\newpage
\section{A stronger norm criterion and the certificate criterion}
\begin{proposition}\label{prop:sqrt}
Fix $q$, write $B=L^{-1}$, $N=U-I$, $r=\norm{N}>0$, and define
\[
 K=L\left(I-\frac{NN^T}{r^2}\right)L^T\succeq0.
\]
A sufficient condition for $\norm{(I+\mu B)^{-1}N}\le r$ for every real
$\mu\ge0$ is the single inequality
\begin{equation}\label{eq:sqrtcriterion}
 L+L^T+2K^{1/2}\succeq0,
\end{equation}
where $K^{1/2}$ is the symmetric positive-semidefinite square root.
\end{proposition}
\begin{proof}
Put $P=I+\mu B$. The desired norm inequality is equivalent, by congruence,
to $r^2PP^T-NN^T\succeq0$. Congruence by $L$ and division by $r^2$
turn this into
\[
 Q(\mu)=K+\mu(L+L^T)+\mu^2 I\succeq0.
\]
The exact identity
\[
 Q(\mu)=(\mu I-K^{1/2})^2+
             \mu\bigl(L+L^T+2K^{1/2}\bigr)
\]
proves the assertion. The first term is the square of a symmetric matrix;
the second is positive semidefinite under \eqref{eq:sqrtcriterion}.
\end{proof}
This is a \emph{sufficient} test. It has not been proved for all Radau
stage counts. Failure of this stronger test would not by itself refute
IE-27. In particular, the orientation $PP^T$ and $NN^T$ in this argument
must not be interchanged with $P^TP$ and $N^TN$.

\begin{proposition}[Common energy criterion]\label{prop:energy}
Suppose $H=H^T\succ0$ and $t>0$ satisfy
\[
 HB+B^TH\succeq0,\qquad tH-N^THN\succeq0.
\]
Then $\rho((I+\mu B)^{-1}N)\le\sqrt t$ for every real $\mu\ge0$.
\end{proposition}
\begin{proof}
For $\|x\|_H=(x^*Hx)^{1/2}$,
\[
 (I+\mu B)^TH(I+\mu B)
 =H+\mu(HB+B^TH)+\mu^2B^THB\succeq H.
\]
Thus the inverse is nonexpansive in this norm, while
$\|Nx\|_H\le\sqrt t\,\|x\|_H$. The product has induced
$H$-norm at most $\sqrt t$, which also bounds its spectral radius.
\end{proof}
The exact certificates use Proposition \ref{prop:energy}; the numerical
square-root experiment is \emph{not} their proof mechanism.

\newpage
\section{Two additional exact certificates}
The new files \code{new_certificates/q096.json} and
\code{new_certificates/q128.json} pass the retained integer-interval
verifier. They contain rational $H,t,v$ and rational isolating intervals
for every interior Radau node. The added squared radii are
\[
 t_{96}=\frac{7801459205}{10^{10}},\qquad
 t_{128}=\frac{8069729082}{10^{10}}.
\]
For each of these stages, the verifier establishes
\[
 H\succ0,\quad HB+B^TH\succ0,\quad tH-N^THN\succ0,
 \quad I-N^TN\succ0,
\]
\[
 v\ne0,\qquad \norm{Nv}^{\,2}-t\norm{v}^{\,2}>0.
\]
The last two checks prove $\sqrt t<r<1$. Combining them with
Proposition \ref{prop:energy} gives the \emph{continuum} statement
\[
 \rho((I+\mu B_q)^{-1}N_q)\le\sqrt{t_q}<r_q<1
 \qquad\text{for every real }\mu>0,
 \quad q\in\{96,128\}.
\]

\begin{center}\small
\begin{tabular}{@{}lrr@{}}
\toprule
Verified positive lower bound & $q=96$ & $q=128$\\\midrule
Smallest $L$ pivot & $70.629700811196$ & $94.174566164618$\\
Smallest LDL pivot of $H$ & $1.001307579999$ & $1.002030669999$\\
Smallest LDL pivot of $HB+B^TH$ & $0.000156834663$ & $0.000110365601$\\
Smallest LDL pivot of $tH-N^THN$ & $0.376277528400$ & $0.436040730349$\\
Smallest LDL pivot of $I-N^TN$ & $0.716423169363$ & $0.716128872382$\\
Rayleigh gap for supplied $v$ & $0.001562635417$ & $0.001616369979$\\\bottomrule
\end{tabular}
\end{center}
These are downward-rounded displays of certified rational lower bounds,
not numerical eigenvalue estimates. LDL pivots are not eigenvalues.

\subsection*{Trust boundary and finite scope}
Node brackets use 256 bits; all interval operations use integer endpoints
with denominator $2^{1024}$ and outward rounding. Polynomial signs isolate
the nodes; exact matrix formulas and triangular elimination enclose the
factors; interval LDL elimination proves positivity. No floating-point
optimizer or eigensolver is trusted by the verifier.

The complete supplied stage set is
\[
 \mathcal S=\{2,3,\ldots,64\}\cup\{80,96,128\},\qquad |\mathcal S|=66.
\]
This does \emph{not} mean all stages through 128. In particular, stages
65--79, 81--95, 97--127, and stages above 128 are not covered by these
certificates. Existence of a certificate for every stage has not been proved.

From the extracted archive, \code{python code/run_all.py} runs the
standard-library proof checks. Optional numerical generation is separate.

\newpage
\section{Why two tempting generic arguments are invalid}
\subsection*{6.1 Positive pivots and weighted accretivity are insufficient}
Consider the exact rational matrices
\[
 B=\begin{pmatrix}1&0\\5&1\end{pmatrix},\quad
 L=B^{-1}=\begin{pmatrix}1&0\\-5&1\end{pmatrix},\quad
 N=\begin{pmatrix}0&1/10\\0&0\end{pmatrix}.
\]
They have the stated triangular forms, positive pivots, and $\norm{N}=1/10<1$.
Moreover,
\[
 D=L(I+N)=\begin{pmatrix}1&1/10\\-5&1/2\end{pmatrix},\qquad
 H_0=\operatorname{diag}(50,1),
\]
\[
 H_0D+D^TH_0=\operatorname{diag}(100,1)\succ0.
\]
Nevertheless, at the positive real shift $\mu=1$,
\[
 (I+B)^{-1}N=\begin{pmatrix}0&1/20\\0&-1/8\end{pmatrix},
 \qquad \rho((I+B)^{-1}N)=\frac18>\frac1{10}=\norm{N}.
\]
The script \code{code/exact_obstructions.py} checks these identities
and inequalities using rational arithmetic.

\textbf{This is not a Radau matrix and is not a counterexample to IE-27.}
It proves that positive pivots, triangularity, and diagonal weighted
accretivity of $D$ do not suffice for a general argument. Any proof that
uses only those properties has discarded essential structure.

\subsection*{6.2 The Euclidean inverse factor is not always nonexpansive}
The recovered exact three-stage matrix is
\[
 B_3=\begin{pmatrix}
 \frac45-\frac{\sqrt6}{5}&0&0\\
 \frac9{50}+\frac{29\sqrt6}{200}&\frac3{10}+\frac{3\sqrt6}{40}&0\\
 \frac49-\frac{\sqrt6}{36}&\frac49+\frac{\sqrt6}{36}&\frac19
 \end{pmatrix}.
\]
For $w=(0,1,-2)^T$,
\[
 w^TB_3w=\frac{-52+7\sqrt6}{360}<0.
\]
Hence
\[
 \norm{(I+\mu B_3)w}^{\,2}-\norm{w}^{\,2}
 =2\mu w^TB_3w+\mu^2\norm{B_3w}^{\,2}<0
\]
for sufficiently small real $\mu>0$. The inverse factor then has
Euclidean norm greater than one. This invalidates the shortcut
$\norm{(I+\mu B)^{-1}}\le1$, but does not invalidate the spectral bound
for the product with $N$.

\newpage
\section{An exact failure for complex right-half-plane shifts}
There is also a genuine Radau example outside the stated shift domain.
For $q=3$, take
\[
 \mu=\frac1{100}+3i,\qquad R_0=\frac{41}{100}.
\]
Exact arithmetic in $\mathbb Q(\sqrt6,i)$ establishes
\begin{equation}\label{eq:complex}
 \rho((I+\mu B_3)^{-1}N_3)>R_0>\norm{N_3}.
\end{equation}
Here $\Re\mu>0$, but $\mu$ is \emph{not real}. Thus
\eqref{eq:complex} does not refute IE-27; it prevents an unjustified
extension of the target to the entire right half-plane.

\subsection*{Exact proof certificate}
Let $s=\sqrt6$. Together with the $B_3$ on the previous page, use
\[
 N_3=\begin{pmatrix}
 0&-\frac{53}{25}+\frac{76s}{75}&\frac{16}{25}-\frac{22s}{75}\\
 0&0&\frac2{25}+\frac{3s}{25}\\
 0&0&0
 \end{pmatrix}.
\]
The checker first proves $R_0^2I-N_3^TN_3\succ0$ by exact leading
principal minors, establishing the second strict inequality in
\eqref{eq:complex}.

Set $Y=(I+\mu B_3)^{-1}N_3$. Its first column is zero. Its two nontrivial
eigenvalues are the roots of $z^2+az+b$, with
\[
 a=-Y_{22}-Y_{33},\qquad b=Y_{22}Y_{33}-Y_{23}Y_{32}.
\]
Put $A=a/R_0$ and $E=b/R_0^2$. If both roots had modulus at most $R_0$,
then necessarily
\begin{equation}\label{eq:schur}
 |A-\overline A E|\le1-|E|^2.
\end{equation}
Indeed, writing the scaled roots as $z_1,z_2$, with moduli $u,v\le1$,
\[
 A-\overline A E=-z_1(1-|z_2|^2)-z_2(1-|z_1|^2),
\]
whose modulus is at most $(u+v)(1-uv)\le(1+uv)(1-uv)=1-|E|^2$.

The supplied checker obtains the two strict inequalities
\[
 1-|E|^2>0,\qquad
 |A-\overline A E|^2-(1-|E|^2)^2>\frac{27}{100}.
\]
They contradict \eqref{eq:schur}. Every field operation is exact; signs
are bounded using the checked rational bracket
$2449489/10^6<\sqrt6<2449490/10^6$. Full rational components and bounds
are in \code{results/exact_obstructions.json}; no numerical root-finding
is part of this proof.

\newpage
\section{Numerical observations, reproducibility, and the gap}
The optional script \code{code/diagnostics.py} evaluates the sufficient
matrix inequality \eqref{eq:sqrtcriterion} and samples the shift family.
Its output is explicitly marked \code{NUMERICAL_ONLY}. Selected estimates
of the smallest eigenvalue of $L+L^T+2K^{1/2}$ are
\begin{center}
\begin{tabular}{@{}rrrrrr@{}}
\toprule
$q$&2&3&16&64&128\\\midrule
Numerical minimum&1.5943&0.7016&2.0536&2.1049&2.1238\\\bottomrule
\end{tabular}
\end{center}
The exact $K$ is singular positive semidefinite; the numerical script
clips negative roundoff before taking its square root. These numbers are
neither exact positivity certificates nor an argument for untested stages.
Likewise, maxima over the recorded 101-point shift grid are not maxima
over all positive shifts. No conclusion depends on extrapolation from them.

\subsection*{Remaining mathematical obligation}
The unrestricted statement \eqref{eq:target} remains to be proved for
$q\notin\mathcal S$ in the intermediate shift regime of Section 3, or
refuted by a counterexample satisfying \emph{all} of its hypotheses.
One sufficient route would be a proof of \eqref{eq:sqrtcriterion} for
all stages. Another would be a construction of the common-energy
certificates for every stage with $t\le r_q^2$. Neither existence result
is established here. The finite certificate generator is not known to
succeed for every stage and is not presented as such a proof.

\subsection*{Files and independent checks}
The top-level report is this continuation; \code{prior_results/report.pdf}
is the recovered earlier manuscript. The new exact checks are in
\code{code/}, the added certificates in \code{new_certificates/}, and
machine-readable outcomes in \code{results/}. Run
\begin{quote}
\code{python code/run_all.py}
\end{quote}
for all 66 supplied certificates, the exact low-stage checks, and the
regression tests. Run \code{python code/exact_obstructions.py} to check
only the two rejected extensions. Verification requires only the Python
standard library. Optional candidate generation and numerical diagnostics
require NumPy and SciPy. The SHA-256 manifest checks file integrity, not
mathematical truth.

\vspace{2pt}
\begin{thebibliography}{9}\footnotesize
\bibitem{problem} \emph{OpenProblemsInNLA} repository, IE-27, ``Spectral disk for a Radau stage preconditioner.''
Problem statement checked 12 September 2026; page status dated
11 September 2026.\newline
\url{https://github.com/ajt60gaibb/OpenProblemsInNLA/tree/main/linear-systems-and-elimination/IE-27}
\bibitem{outrata} M. Outrata, \emph{On the recent advances of spectral
analysis for systems arising from fully-implicit RK methods},
arXiv:2510.21241v1 (2025), Sections 2.1--2.3.\newline
\url{https://arxiv.org/html/2510.21241v1}
\end{thebibliography}
\end{document}
